% Transcription — fonds Alexandre Grothendieck, shelfmark 31, batch 2 (pages 21-40).
% Established from the facsimile archives/batches/31/batch-02.pdf (a mirror of
% https://grothendieck.umontpellier.fr/31.pdf), per the skill
% transcribe-grothendieck.
%
% Pass: Fable 5.1 (claude-fable-5-1), 2026-09-10 — first pass, unchecked against
% the pages by a human.
%
% The batch closes the run begun on page 19 (batch 1), then holds two runs
% opened by cover sheets inscribed in his hand (pages 23, 32), which become
% the section headings and carry no \page{}.
%
%   Pages 21-22 — ink, fair: the end of « Théorèmes de M. Artin ». Corollaire
%     2 (K of finite type over R, X a model), Corollaire 3 (X of finite type
%     over R: X(R) = ∅ ⟺ cd_2(X) < ∞, and then H^i(X,F) ≃ H^i(X(C)/Z/2Z, F)
%     with cd_2(X) ≤ 2 dim X), Théorème 2 (k a field, r a prime: the four
%     equivalent conditions on cd_r(k') = cd_r(k) + deg tr k'/k, with the
%     margin list of the cases where they hold), and a Corollaire on a local
%     noetherian domain A: cd_r(K) ≥ cd_r(k) + dim A, with the two possible
%     exceptions (r = car K, or r = 2 and an orderable residue field k(x)
%     with cd_2 k(x)(√−1) finite), the N.B. on the catenary universally
%     japanese case, and the margin's « cette restriction est-elle
%     essentielle ??? ».
%   Pages 24, 26, 28, 30 — ink, on the backs of four typescript leaves: the
%     « Lemme de Cartan ». Page 24 states it — T a site, F an abelian sheaf,
%     n ∈ Z, M a class of objects stable under fibre products whose coverings
%     are cofinal — with the three equivalent conditions (i) H^i(S,F) = 0,
%     (ii) Ȟ^i(U,F) = 0 for every covering U by objects of M, (iii)
%     lim_U Ȟ^i(U,F) = 0, all for 0 < i ≤ n. Page 26 proves it through the
%     spectral sequence Ȟ^p(U, H^q(F)) ⟹ H^*(S,F), by induction on n, with
%     a sketch of the (p,q)-plane in the margin. Page 28 is a second Lemme
%     with the same notations: (i) lim Ȟ^i(U,F) ≃ H^i(S,F) for i ≤ n, (ii)
%     and (iii) the vanishing of lim Ȟ^p(U, H^q(F)) and of lim Ȟ^p(U,
%     H̄^q(F)), where H̄^q(F)(T) = lim_V Ȟ^q(V,F), and the cochain-level
%     conditions (iii bis), (iii ter); page 30 proves the implications
%     (« Cas 2 »), again by induction on n, and ends with a diagram of what
%     is established and a « Solution » naming the four implications wanted.
%   Pages 27, 25 — two leaves of a typescript on formal completion along a
%     closed subscheme (Théorème 2.1 on the S-module K^i = R^i f_*(gr_J F)
%     and the Mittag-Leffler condition, Théorème 2.2 on an adic noetherian
%     ring with an ideal of definition generated by one t; then Remarques
%     5.11 on the comparison of étale coverings and invertible modules of
%     X' and X̂', and Picard schemes of non-proper schemes after Murre),
%     corrected in his hand — insertions in \add{}, the hats the machine did
%     not carry restored where inked — and given in text order, which is the
%     reverse of the archive's.
%   Pages 34-38 — blue ink, a fair hand which is not his, on letterhead of
%     the École normale supérieure, the first leaf headed « GRÖTHENDIECK. »:
%     Th 1, for a discrete group G acting simplicially on a connected
%     simplicial set E with quotient B = E/G, im(π_1(E) → π_1(B)) is normal
%     and π_1(B)/im p_* = G/S, S the normal subgroup generated by the
%     stabilisers of the vertices, E/S the corresponding covering; the
%     definitions of edges, paths, loops and homotopy; Lemme 1 (paths lift),
%     Lemme 2 (the image is normal), the map χ: G → π_1(B)/im π_1(E), its
%     surjectivity, multiplicativity and kernel; Lemme 3 (E/S is a normal
%     covering of E/G with group G/S); Th 2 (G finite of order n: the image
%     of H_*(E) → H_*(B) contains the elements divisible by n); Th 3 (E
%     contractible, S finite of order n: π_i(B) has all elements of order
%     dividing a power of n for i > 1), whose proof breaks off.
%
% Pages 29 and 31 are typescript leaves of the same text (Corollaire 2.2,
% Lemme 2.3, Corollaire 2.4) carrying nothing in his hand — absent. Pages
% 33, 37, 39 and 40 are blank — absent.
%
% The dating is the inventory's own, for the whole shelfmark.
\documentclass[11pt,a4paper]{article}
% Le préambule commun à toutes les transcriptions du fonds Grothendieck.
%
% Il tient en une idée : **l'incertitude se compose, elle ne se résout pas.**
% Une transcription qui lisse un mot douteux a détruit la seule chose pour
% laquelle on la faisait — un lecteur doit pouvoir savoir, à chaque endroit,
% ce qui était sur le feuillet et ce qui a été ajouté. Les six macros qui
% suivent sont donc l'appareil critique, et non de la décoration.
%
% Les mêmes commandes sont reconnues par `scripts/render.mjs`, qui produit la
% vue de lecture du site. Une macro ajoutée ici doit l'être aussi là-bas,
% faute de quoi le rendu échouera bruyamment — ce qui est voulu : un rendu
% silencieusement faux est pire qu'un rendu absent.

\ProvidesPackage{grothendieck}[2026/08/08 Transcriptions du fonds Grothendieck]

\usepackage{amsmath,amssymb,amsthm}
\usepackage{mathrsfs}
\usepackage{stmaryrd}
% Les diagrammes commutatifs sont partout dans ce fonds : c'est la langue
% même de la géométrie algébrique de Grothendieck, et une transcription qui
% les rendrait en prose ne transcrirait rien.
\usepackage{tikz-cd}
% Français par défaut — c'est la langue des feuillets — mais l'anglais est
% chargé aussi : la traduction et le résumé sont des documents anglais, et la
% typographie française (espaces avant les deux-points) y serait une faute.
% Ces documents basculent par \selectlanguage{english} en tête de corps.
\usepackage[english,french]{babel}
\usepackage{fontspec}
\usepackage{geometry}
\geometry{a4paper,margin=25mm,marginparwidth=28mm}
\usepackage{marginnote}
\usepackage{xcolor}
% ulem, not soul. `soul`'s \st and \ul are built on a character-by-character
% scan that XeLaTeX's font machinery breaks: the document fails with an
% undefined control sequence rather than degrading. ulem does the same two jobs
% and is Unicode-engine safe. `normalem` stops it hijacking \emph, which the
% transcriptions use for Grothendieck's own emphasis.
\usepackage[normalem]{ulem}
\usepackage{hyperref}
\hypersetup{colorlinks=true,linkcolor=black,urlcolor=black,pdfborder={0 0 0}}

% --- Métadonnées du lot -----------------------------------------------------
% Reprises telles quelles par le script de rendu, qui les affiche en tête de
% la vue de lecture. Elles fixent ce que le fichier prétend couvrir : sans
% elles, une transcription flotte sans point d'attache dans un fonds de seize
% mille feuillets.

\newcommand{\folder}[1]{\gdef\grothFolder{#1}}
\newcommand{\batch}[1]{\gdef\grothBatch{#1}}
\newcommand{\leaves}[2]{\gdef\grothFirst{#1}\gdef\grothLast{#2}}
\newcommand{\foldertitle}[1]{\gdef\grothFoldertitle{#1}}
\gdef\grothFolder{?}\gdef\grothBatch{?}\gdef\grothFirst{?}\gdef\grothLast{?}\gdef\grothFoldertitle{}

% --- Le résumé --------------------------------------------------------------
%
% Il ouvre la lecture modernisée, sous le titre « Résumé », et il est composé
% en petit. Ce n'est pas un ornement : c'est le seul passage du document qui
% ne soit pas des mathématiques, et un lecteur doit pouvoir le reconnaître
% avant de l'avoir lu — puis le sauter s'il connaît déjà le sujet, ou s'y
% arrêter s'il ne le connaît pas. Le filet qui le suit marque la reprise du
% corps.

\newenvironment{resume}
  {\small\setlength{\parskip}{0.45em}}
  {\par\medskip\hrule\medskip}

% --- Mots-clés ---------------------------------------------------------------
%
% En anglais, en clôture du résumé de la lecture modernisée : le vocabulaire
% moderne sous lequel chercher ce que les feuillets construisent. Le manifeste
% du site (`npm run manifest`) les extrait pour étiqueter la cote entière —
% cette ligne est la source unique des tags, il n'y a pas de fichier à côté.
\newcommand{\keywords}[1]{\par\smallskip\noindent
  {\footnotesize\textsc{Keywords} --- \textit{#1}}\par}

% --- L'appareil critique ----------------------------------------------------

% Le numéro de feuillet, en marge. C'est l'ancre qui permet de régler un
% désaccord sur une formule en regardant *un* feuillet plutôt qu'une chemise
% de six cents. Le site s'en sert aussi pour tourner les pages du fac-similé
% quand on fait défiler la transcription.
\newcommand{\leaf}[1]{%
  \marginnote{\footnotesize\color{gray!70}\textsf{#1}}%
  \label{leaf:#1}\ignorespaces}

% Illisible : on ne devine pas. Un mot inventé qui se lit comme les autres est
% le pire résultat possible.
\newcommand{\ill}{\textcolor{red!60!black}{[\,\ldots\,]}}

% Lecture douteuse : le mot est donné, et signalé comme incertain.
\newcommand{\uncertain}[1]{\uline{#1}}

% Ajout de l'éditeur — une lettre manquante, un mot sous-entendu.
\newcommand{\add}[1]{\textcolor{blue!50!black}{[#1]}}

% Biffé par Grothendieck lui-même. Ce qu'il a rayé fait partie du document :
% une rature dit souvent où la pensée a changé de direction.
\newcommand{\struck}[1]{\sout{#1}}

% Note du transcripteur, sur l'état matériel du feuillet.
\newcommand{\note}[1]{\par\smallskip{\small\color{gray!60!black}\textit{[#1]}}\par\smallskip}

% Note marginale de Grothendieck, distincte du corps du texte.
\newcommand{\marginal}[1]{\par\smallskip\noindent\hspace{1em}%
  {\small\color{gray!60!black}#1}\par\smallskip}

% --- Titre ------------------------------------------------------------------

\AtBeginDocument{%
  \begin{center}
    {\large\bfseries Fonds Alexandre Grothendieck --- cote n\textsuperscript{o}~\grothFolder}\\[2pt]
    {\itshape\grothFoldertitle}\\[4pt]
    {\small feuillets \grothFirst--\grothLast\ (lot \grothBatch)}\\[2pt]
    {\footnotesize\color{gray!60!black}Transcription établie d'après le fac-similé numérisé
     par l'Université de Montpellier.\\
     Les lectures incertaines sont soulignées, les illisibles notées [\,\ldots\,],
     les ajouts entre crochets.}
  \end{center}
  \vspace{1em}}

\endinput


\folder{31}
\batch{2}
\pages{21}{40}
\dating{[1963-1973]}
\watermark{Édition de démonstration}
\foldertitle{SGA A [SGA 4] (Résidus de rédaction) : notes manuscrites (s.d.)}

\begin{document}

\section*{Théorèmes de M. Artin, suite (pages 21 à 22)}

\note{Suite de la page 19 (lot 1), dont le Théorème 1 et le Corollaire 1
fixent les notations : $K$ extension de type fini de $\mathbb{Q}$, $X$ un
modèle, les conditions (i) à (iii bis).}

\page{21}
\textbf{Corollaire 2} Soit $K$ extension de type fini de $\mathbb{R}$, $X$
un modèle. Conditions équivalentes (i) (i bis) (ii) (ii bis) (ii ter)
(ii quater) (iii) (iii bis).
\struck{(i) $-1$ une somme de carrés dans $K$}

\textbf{Corollaire 3} $X$ préschéma de type fini sur $\mathbb{R}$.
Conditions équivalentes
\begin{enumerate}
\item[(i)] $X(\mathbb{R}) = \emptyset$
\item[(ii)] $cd_2(X) < +\infty$
\end{enumerate}
\struck{(iii)} Lorsqu'il en est ainsi, alors pour tt faisceau (étale) $F$
sur $X$, on a
\[
H^i(X,F) \simeq H^i(X(\mathbb{C})/\mathbb{Z}/2\mathbb{Z},\, F)
\]
\note{Sous le second membre, une accolade porte « topologie ordinaire ».}
et \ill{}
\[
cd_2(X) \leqslant 2\dim X
\]

\textbf{Théorème 2} Soit $k$ un corps, \struck{$k$ un \ill{}} \add{$r$ un
\struck{\ill{}} \uncertain{nombre premier}}. Conditions équivalentes
\note{L'insertion « $r$ un … premier » est cerclée et rattachée à la ligne
par un trait.}
\begin{enumerate}
\item[(i)] Pour tt extension de type fini $k'$ de $k$, on a
\[
cd_r(k') = cd_r(k) + \deg.\ \mathrm{tr.}\ k'/k
\]
\item[(i bis)] Comme dessus, \struck{\ill{} avec $k'$ une extension
quadratique} $k' = k(\sqrt{-1})$ ! \struck{et même $k' = k(\sqrt{-1})$}
\item[(ii)] $k'$ contient un corps $k_0$ de type fini sur le corps premier,
et tel que $cd_r(k_0) < +\infty$, ou $cd_2\, k(\sqrt{-1}) = +\infty$
\note{L'apostrophe de $k'$ est sur la page ; la condition semble porter
sur $k$.}
\item[(iii)] On n'est pas dans le cas suivant : $k$ est ordonnable
(\struck{donc} i.e.\ de car $0$ [et $-1$ n'est pas somme de carrés]),
$r = 2$ [$cd_2(k(\sqrt{-1})) < +\infty$. [d'où $cd_2(k) = +\infty$]
\end{enumerate}
\note{Les étiquettes (i) et (iii) sont cerclées.}

\marginal{Conditions tjs \uncertain{satisfaites} \ill{} si l'une des
conditions suiv.\ : $\operatorname{car} k > 0$ ; $r \neq 2$ ; $-1$ somme de
carrés ; \struck{\ill{}} $cd_2\, k(\sqrt{-1}) < +\infty$ ;
$[cd_r(k) < +\infty]$}
\note{Note marginale écrite en oblique, de bas en haut, à gauche du
Théorème 2.}

\page{22}
\textbf{Corollaire} $A$ anneau local \struck{hensélien} noethérien intègre,
corps des fractions $K$, corps résiduel $k$. Alors pour tt premier $r$, on a
\[
cd_r(K) \geqslant cd_r(k) + \dim A
\]
\struck{et $A$ univ.\ japonais / univ.\ caténaire} \struck{[avec égalité
dans le cas $A$ hensélien ?]} \uncertain{: la} seule exception possible
dans les cas suivants :
\begin{enumerate}
\item[a)] $r$ est la caractéristique de $K$, (\uncertain{i.e.}\
caractéristique $> 0$)
\item[b)] $r = 2$, \struck{$A$ un anneau de dim $> 0$, il existe}
\struck{$A$ « d'inégales caractéristiques » ou la caractéristique
résiduelle est $2$, \uncertain{ou} un idéal premier $\mathfrak{p} \subset
A$, \ill{} avec $A/\mathfrak{p}$ de dim $1$,} $\exists\, x \in
\operatorname{Spec}(A)$ tel que $k(x)$ soit ordonnable, et $cd_2\,
k(x)(\sqrt{-1})$ soit fini (corps réel \ill{})
\end{enumerate}
\note{Dans b), le passage biffé est écrit sur trois lignes et rayé d'un
trait continu ; l'accolade de gauche réunit a) et b).}

\marginal{Ex.\ $A = \mathbb{R}[X,Y]/(X^2+Y^2)$ \uncertain{localisé}}

N.B.\ Si $A$ est caténaire et univ.\ japonais, (ce qui garantit le fait que
les lieux normaux des \ill{} de $A$ sont ouverts) \add{dans le cas b)},
alors \struck{il suffit} on peut préciser \struck{\ill{}} l'exception
possible \uncertain{par} : $k$ lui-même « corps réel \ill{} ». Enfin,
lorsque $A$ est régulier, cette exception elle-même est inoffensive, i.e.\
on a l'inégalité annoncée, (pouvant \ill{} dans le cas \ill{})].
\note{« dans le cas b) » est une insertion interlinéaire, cerclée.}

\marginal{cette restriction est-elle essentielle ???}

\section*{Comparaison de topologies --- Lemme de Cartan (pages 24 à 30)}

\note{Le titre est de la main de Grothendieck, sur la chemise qui précède
(page 23) : « Comparaison de topologies — Lemme de Cartan ». Dans ce
feuillet, $\mathcal{M}$ est une classe d'objets du site $T$, $\mathcal{U}$,
$\mathcal{V}$ des recouvrements, $\mathcal{H}^q(F)$ le préfaisceau $U
\mapsto H^q(U,F)$ (écrit $H^q(F)$ sur la page), et $\overline{H}^q(F)$ le
préfaisceau défini page 28.}

\page{24}
\textbf{Lemme de Cartan}

$T$ un site, \struck{$S \in \operatorname{Ob} T$}

$F$ faisceau abélien sur $T$

$n \in \mathbb{Z}$

\struck{$\mathcal{U}$ (familles \uncertain{filtrantes} croissantes de
rec.\ (de $S$) \ill{}}

$\mathcal{M}$ \uncertain{sous-ensemble} des $\operatorname{Ob} T$
\struck{formé de} --- a) $\mathcal{M}$ stable par produits fibrés ;
b) pour tt $S \in \mathcal{M}$, \struck{\ill{}} les \emph{recouvrements}
de $S$ formés par des $\in \mathcal{M}$ sont cofinaux parmi les rec.\ [il
suffit \ill{} $\subset$, qu'ils suffisent : \ill{} \struck{$H(S)$} tt
\uncertain{calcul} \ill{} $H^q(S,F)$, $0 \leqslant q \leqslant n$.]
\note{Les étiquettes a) et b) sont cerclées ; le crochet est rattaché par
un trait à « cofinaux ».}

Conditions équivalentes
\begin{enumerate}
\item[(i)] $\forall S \in \mathcal{M}$, $H^i(S,F) = 0$ pour $0 < i
\leqslant n$
\item[(ii)] $\forall S \in \mathcal{M}$, \struck{\ill{}} et tt rec.\
\add{$\mathcal{U}$} de $S$ [formé par des $S_i \in \mathcal{M}$, \ill{}]
\[
\check{H}^i(\mathcal{U},F) = 0 \quad \text{pour } 0 < i \leqslant n
\]
\item[(iii)] $\forall S \in \mathcal{M}$, \ill{}
\[
\varinjlim_{\mathcal{U}} \check{H}^i(\mathcal{U},F) = 0 \quad \text{pour }
0 < i \leqslant n
\]
où $\mathcal{U}$ varie dans les rec.\ de $S$ formés d'objets $\in
\mathcal{M}$
\end{enumerate}
\note{Dans (ii), le $\mathcal{U}$ est cerclé et inséré au-dessus de la
ligne.}

\page{26}
Cas \ill{}. (i) $\Longrightarrow$ (ii) par suite spectrale
\note{Le numéro du cas est réduit à un point ; « Cas 2 » ouvre la page 30.}
\[
H^{*}(S,F) \Longleftarrow \check{H}^p(\mathcal{U}, \mathcal{H}^q(F)) =
E_2^{pq}
\]
l'hyp.\ (i) implique $E_2^{pq} = 0$ pour $0 < q \leqslant n$, et
\struck{l'hyp.\ \ill{} pour \struck{$E_2^{pq} = 0$ pour $q = 0$}}
d'où
\[
H^i(S,F) \overset{\simeq}{\Longleftarrow} \check{H}^i(\mathcal{U},F)
\quad \text{pour } i \leqslant n
\]
et par suite $\check{H}^i(\mathcal{U},F) = 0$ pour $0 < i \leqslant n$.

(ii) $\Longrightarrow$ (iii) trivial

(iii) $\Longrightarrow$ (i) par récurrence sur $n$. C'est \struck{trivial}
si $n \leqslant 0$. Supposons $n > 0$, et l'assertion prouvée pour $n-1$.
\struck{Alors \ill{} Regardons} \struck{$E_2^{pq} = 0$ si $0 < q < n$}
\[
H^{*}(S,F) \Longleftarrow \varinjlim_{\mathcal{U}} \check{H}^p(\mathcal{U},
\mathcal{H}^q(F)) = E_2^{pq}
\]
Par hyp.\ de récurrence, on a
\begin{enumerate}
\item[a)] $E_2^{pq} = 0$ si $0 < q < n$
\end{enumerate}
\struck{de plus, $E_2^{pq}$ \ill{}}
Par hyp.\ (iii), on a
\begin{enumerate}
\item[b)] $E_2^{pq} = 0$ si $q = 0$, $0 < p \leqslant n$
\end{enumerate}
Enfin on a \struck{d'après l'hyp.\ \uncertain{(A)} \ill{}}
\begin{enumerate}
\item[c)] $E_2^{pq} = 0$ si $p = 0$, $q > 0$
\end{enumerate}
On trouve donc $E_2^{pq} = 0$ si $p + q = n$, d'où $H^n(S,F) = 0$, cqfd.
\note{Un croquis dans la marge gauche figure le plan $(p,q)$ : des croix
sur les lignes horizontales $q \geqslant n$, des croix sur l'axe des $p$ à
partir de $p = n$ (le point $n$ et « $n+1$ » y sont marqués), et la
diagonale $p + q = n$ en pointillé, qui ne rencontre aucune croix.}

\page{28}
\textbf{Lemme.} Notations $T$, $F$, $n$, $\mathcal{M}$ comme dans le lemme
de Cartan. \struck{Conditions équivalentes} \add{Considérons les}
\note{« Considérons les » est écrit au-dessus de « Conditions
équivalentes » biffé.}
\begin{enumerate}
\item[(i)] $\forall S \in \mathcal{M}$,
\[
\varinjlim_{\mathcal{U}} \check{H}^i(\mathcal{U},F) \simeq H^i(S,F), \qquad
i \leqslant n
\]
[$\mathcal{U}$ parcourt les rec.\ de $S$ formés de $S_i \in \mathcal{M}$,]
\item[(ii)] $\forall S \in \mathcal{M}$, \struck{et tt rec.\ $\mathcal{U}$
de $S$ formé d'éléments de $S$, on ait $H$ élément :}
\[
\varinjlim_{\mathcal{U}} \check{H}^p(\mathcal{U}; \mathcal{H}^q(F)) = 0
\quad \text{pour } \uncertain{p < q \leqslant n}
\]
\struck{$0 < q \leqslant n$} \struck{et $q \neq 0$}
\note{Le $0$ de « $0 < q \leqslant n$ » est repassé en $p$ et le signe qui
suit est resté un $<$ ; la ligne biffée de (iii bis) écrit $p + q
\leqslant n - 1$. Les deux conditions biffées sont écrites sous la
première.}
\item[(iii)] $\forall S \in \mathcal{M}$,
\[
\varinjlim_{\mathcal{U}} \check{H}^p(\mathcal{U}, \overline{H}^q(F)) = 0
\quad \text{pour } \uncertain{p < q \leqslant n} \quad \struck{\text{et }
q \neq 0}
\]
où on pose ($\forall T \in \mathcal{M}$)
\[
\overline{H}^q(F)(T) = \varinjlim_{\mathcal{V}} \check{H}^q(\mathcal{V},
\uncertain{F})
\]
\note{Le second argument se lit plutôt $T$ sur la page.}
$\mathcal{V}$ parcourant les rec.\ de $T$ formés d'objets $\in
\mathcal{M}$.
\item[(iii bis)] $\forall S \in \mathcal{M}$, $\mathcal{U}$ rec.\ de $S$
par des $U_i \in \mathcal{M}$, et $\xi \in C^p(\mathcal{U}, \mathcal{H}^q(F))$
\struck{($p + q \leqslant n-1$, $q \neq 0$)} $0 < q \leqslant n$, il existe
$\mathcal{U}'$ rec.\ de $S$ \uncertain{fini} par $U'_i \in \mathcal{M}$,
majorant $\mathcal{U}$, \struck{\ill{}} et $\mathcal{U}' \to \mathcal{U}$,
tel que l'image de $\xi$ dans $C^p(\mathcal{U}', \mathcal{H}^q(F))$ soit
nulle.
\item[(iii ter)] \struck{\ill{}} Comme (iii bis), mais $\mathcal{H}^q(F)$
remplacé par $\overline{H}^q(F)$.
\end{enumerate}
\note{Une accolade au crayon réunit (i) et (ii), une autre (iii bis).}

Alors on a \struck{(iii bis) $\Longrightarrow$ (ii), $\Downarrow$}
\[
\begin{array}{ccc}
\text{(ii)} & \Longleftrightarrow & \text{(iii bis)} \\
& & \Downarrow \\
\text{(iii ter)} & \Longrightarrow & \text{(iii)} \\
& & \Downarrow \\
& & \text{(i)}
\end{array}
\]
\note{Un premier diagramme, à gauche, est biffé ; dans celui de droite les
étiquettes de la première colonne sont retouchées, et leur lecture est
douteuse.}

\page{30}
Cas 2. \uncertain{suite} spectrale
\[
H^{*}(S,F) \Longleftarrow \varinjlim_{\mathcal{U}} \check{H}^p(\mathcal{U},
\mathcal{H}^q(F)) = E_2^{pq}.
\]
et
\[
E_2^{pq} = 0 \quad \text{si } p = 0,\ q > 0
\]
Cela \struck{prouve} implique déjà (ii) $\Longrightarrow$ (i),
\uncertain{Or} \uncertain{moyennant} (i), (ii) et (iii) sont équivalents,
\uncertain{donc} (ii) $\Longrightarrow$ (iii). De \uncertain{même}
\uncertain{on a} trivialement (iii bis) $\Longrightarrow$ (ii) \add{donc
(iii bis) $\Longrightarrow$ (i),} et (iii ter) $\Longrightarrow$ (iii),
\struck{et comme (ii) $\Longrightarrow$ (i),} et \struck{que} \ill{}
\uncertain{moyennant} (ii), (iii bis) et (iii ter) sont équivalents,
\ill{} \uncertain{(iii ter)} $\Longrightarrow$ (iii bis). Tout revient
\uncertain{donc} \add{à} \ill{} prouver que (iii) $\Longrightarrow$ (i).
Récurrence sur $n$. Si $n \leqslant 0$, c'est \uncertain{clair}. Supposons
$n > 0$, et que c'est prouvé pour $n - 1$. Donc on sait déjà que
\[
\mathcal{H}^i(F) \overset{\simeq}{\Longleftarrow} \overline{H}^i(F) \quad
\text{pour } i < n \quad \struck{\ill{}}
\]
\struck{\ill{} pour $i = n$.} sur des \uncertain{arguments} $\in
\mathcal{M}$, donc
\[
E_2^{pq} = 0 \quad \text{si } \struck{p + q \leqslant n-1 \text{ et } q
\neq 0}\ 0 < q \leqslant n-1
\]
\struck{Donc} \ill{} la suite spectrale donne \struck{$H^n(F)
\Longrightarrow$}
\[
\varinjlim_{\mathcal{U}} \check{H}^n(\mathcal{U},F) \simeq H^n(S,F)
\]
\struck{\ill{} $H^n(F)$} cqfd.
\note{Un second croquis du plan $(p,q)$ occupe la marge gauche : des croix
sur les deux lignes horizontales supérieures et sur l'axe des $p$ à partir
d'un point marqué « $n-1$ », et deux droites obliques issues de l'axe des
$q$ (marqué « $n$ ») vers ce point.}

\begin{tikzcd}[column sep=small, row sep=small]
\text{(ii)} \arrow[d, Rightarrow] \arrow[dr, Rightarrow] & &
\text{(iii bis)} \arrow[ll, Rightarrow] \arrow[d, Rightarrow] \\
\text{(i)} & \text{(iii)} \arrow[l, Rightarrow] & \text{(iii ter)}
\arrow[l, Rightarrow]
\end{tikzcd}

\note{Diagramme au milieu de la page ; l'étiquette de la flèche du bas,
entre (iii) et (iii ter), est récrite sur un mot gribouillé.}

\[
\begin{array}{ccccc}
\uncertain{\text{(ii bis)}} & \struck{\Longleftrightarrow} & & &
\text{(iii bis)} \\
\Downarrow & & & & \Downarrow \\
\text{(ii)} & \Longrightarrow & \text{(i)} & \Longleftarrow & \text{(iii)}
\end{array}
\]
\note{Diagramme au bas de la page, à gauche ; la double flèche du haut est
biffée. « (ii bis) » n'est pas défini dans le lemme, qui nomme (iii bis)
et (iii ter).}

\textbf{Solution} \uncertain{logique} : Pour avoir les 4 implications
ci-contre, plus le fait que \uncertain{(ii)} $\Longrightarrow$
[\uncertain{(ii)} $\Longrightarrow$ (iii) et (iii bis) $\Longleftrightarrow$
(iii ter)]
\note{Dernières lignes de la page, sous le diagramme ; la première étiquette
entre crochets est écrite « (iii) » et retouchée.}

\subsection*{Deux feuillets de tapuscrit corrigés (pages 27, 25)}

\note{Les pages 27 et 25 sont deux feuillets d'un tapuscrit dont les
versos portent le Lemme de Cartan (pages 26, 24) ; ils sont corrigés de sa
main et donnés dans l'ordre du texte. Les insertions manuscrites sont en
\add{}, les mots biffés en \struck{} ; les accents circonflexes des
complétés ($\widehat{X}$, $\widehat{F}$, $\widehat{U}'$…) et les flèches
que la machine ne portait pas sont restitués là où l'encre les donne.
$\mathcal{J}$ est l'idéal écrit à la main dans les formules (2.5) à (2.8).}

\page{27}
Posons
\begin{gather}
\underline{S} = \operatorname{gr}_{\mathcal{J}} \underline{O}_{X'} =
\coprod_{k \in \underline{N}} \mathcal{J}^k / \mathcal{J}^{k+1} \tag{2.5} \\
\operatorname{gr}_{\mathcal{J}}(\underline{F}) = \coprod_{k \in
\underline{N}} \mathcal{J}^k \underline{F} / \mathcal{J}^{k+1}
\underline{F} \tag{2.6} \\
\underline{K}^i = R^i f_{*}(\operatorname{gr}_{\mathcal{J}}(\underline{F}))
= \coprod_{k \in \underline{N}} R^i f_{*}(\mathcal{J}^k \underline{F} /
\mathcal{J}^{k+1} \underline{F}), \tag{2.7}
\end{gather}
il est clair que $\underline{K}^i$ est muni d'une structure de
$\underline{S}$-Module gradué.

\textbf{Théorème 2.1.} Supposons que $\underline{K}^i$ soit un
$\underline{S}$-Module gradué \emph{de type fini} pour $i = n-1$, $i = n$,
$i = n+1$, alors :
\begin{enumerate}
\item[(i)] $R^n f_{*}(F)$ est cohérent
\item[(ii)] l'homomorphisme $\add{\psi_n}$, (2.3), est un isomorphisme
topologique, \add{la filtration naturelle du deuxième membre de (2.3) est
$\mathcal{J}$-bonne.}
\item[(iii)] le système projectif des $R^n f_{*}(F_k)$ vérifie la
condition de Mittag-Leffler uniforme.
\end{enumerate}
\note{Le $\psi_n$ de (ii) est écrit à la main sur une lettre dactylographiée
illisible ; l'insertion sur la filtration est écrite dans l'interligne.}

La démonstration est très facile à partir de E.G.A.\ $O_{III}$.13.7.7.\
\add{(cf.\ E.G.A.\ III 3.4.2.)}, à condition de rectifier comme suit le
texte page 78 :

\emph{Supprimer}
\begin{quote}
ligne 6 \quad « et $(R^{n+1} T(A_k))_{k \in \underline{Z}}$ »

ligne 13 \quad « et $p + q = n+1$ »

ligne 23 \quad « et pour $n+1$ ».
\end{quote}

\textbf{Théorème 2.2.} Soit $A$ un anneau adique noethérien et soit $I$ un
idéal de définition de $A$. Supposons que $I$ soit engendré par un $t \in
A$. Reprenons les notations (1.31), (1.32), (1.33). Soit $\underline{F}$
un $\underline{O}_{\widehat{X}}$-Module cohérent. Posons
\begin{equation}
F_0 = \underline{F} / \mathcal{J} \underline{F}, \tag{2.8}
\end{equation}
où $\mathcal{J} = t\, \underline{O}_{\widehat{X}}$ est un idéal de
définition de $\widehat{X}$.

\marginal{Soit $T$ une partie fermée de $X' = \operatorname{Spec}(A)$.}
\note{En marge gauche, à la hauteur du Théorème 2.2.}

\page{25}
tel que $A$ soit séparé et complet pour la topologie $J$-adique. Sous ces
conditions :
\begin{enumerate}
\item[(i)] Le foncteur $F \rightsquigarrow \widehat{F}$ de la catégorie des
Modules localement libres sur $U'$ dans la catégorie des Modules localement
libres sur $\widehat{U}'$ est pleinement fidèle.
\item[(ii)] Pour tout Module localement libre $\underline{F}$ sur
$\widehat{U}'$, il existe un Module cohérent $F$ sur $U'$ (pas
nécessairement localement libre, hélas), et un isomorphisme $\widehat{F}
\simeq \underline{F}$.
\end{enumerate}

Il reste seulement à prouver (ii), grâce à 5.9. Or d'après cette remarque,
et exposé IX 2.1.\ il s'ensuit que $\underline{F}$ est induit par un Module
cohérent $\underline{G}$ sur $\widehat{X}'$. D'après le théorème
\struck{de comparaison} d'existence \add{EGA III} 5.1.4., $\underline{G}$
est de la forme $\widehat{F}$, où $F$ est cohérent sur $X$, d'où la
conclusion.

\emph{Remarques} 5.11. \add{1°)} Utilisant 5.10., \add{4.7.} et
\struck{des hypothèses convenables de} une hypothèse convenable, disant que
certains anneaux locaux des fibres géométriques de $X'$ \uncertain{$S'$}
sont « purs » resp.\ parafactoriels, on doit pouvoir obtenir des énoncés
disant \struck{\ill{}} que le foncteur $Z' \rightsquigarrow \widehat{Z}'$
de la catégorie des revêtements étales de $X'$ dans la catégorie des
revêtements étales de $\widehat{X}'$ \struck{\ill{}} (ou ce qui revient au
même, de $X'_0$) est \struck{pleinement fidèle} une équivalence de
catégories, resp.\ que le foncteur $L \rightsquigarrow \widehat{L}$ de la
catégorie des Modules inversibles sur $X'$ dans la catégorie des Modules
inversibles sur $\widehat{X}'$ est une équivalence. Utilisant des résultats
récents de \textsc{Murre}, il est probable que l'on doit pouvoir en déduire
des théorèmes d'existence \struck{pour} des schémas de Picard pour
\add{certains} schémas algébriques \struck{\ill{}} \emph{non propres}.
\struck{\ill{}} De façon générale, l'élimination d'hypothèses de propreté
dans divers théorèmes d'existence, notamment de représentabilité de
foncteurs comme les foncteurs de Hilbert, ou de Picard etc, à l'aide
\note{Le feuillet s'arrête sur « à l'aide » ; la suite n'est pas dans le
lot. Le « $S'$ » après « $X'$ » est dactylographié, sans lien visible.}

\section*{Calculs divers (asphéricité, finitude etc) (pages 34 à 38)}

\note{Le titre est de la main de Grothendieck, au crayon, sur la chemise
qui précède (page 32) : « Calculs divers (asphéricité, \uncertain{finitude}
etc) ». La chemise porte aussi, au crayon : « 4 p.\ Deligne ? + 3 + 6 ».
Les pages 34 à 38 sont d'une écriture qui n'est pas la sienne, à l'encre
bleue, et la page 34 porte en tête « GRÖTHENDIECK. » ; elles sont
transcrites comme le reste. La numérotation « Th 1 », « Th 2 », « Th 3 »,
« Lemme 1 » à « Lemme 3 » est celle de la page.}

\page{34}
\textbf{Th 1} Soit $E$ un ensemble simplicial, $G$ un groupe discret qui
opère sur $E$ \struck{par} par des transformations simpliciales, $B$
\add{$= E/G$} l'ensemble simplicial quotient, \struck{\ill{}} $p : E \to
B$ la projection canonique. Si $E$ est connexe, l'image de $p_{*} : \pi_1(E)
\to \pi_1(B)$ est un sous-groupe distingué de $\pi_1(B)$, et
\[
\pi_1(B) / \operatorname{im} p_{*}(\pi_1(E)) = G/S,
\]
où $S$ est le sous-groupe (distingué) de $G$ engendré par les
stabilisateurs des sommets de $E$, et $E/S$ est le revêtement
\struck{\ill{}} de $B$ correspondant à l'image de $p_{*}(\pi_1(E))$.
\note{« $= E/G$ » est une insertion sous la ligne ; « correspondant à
l'image de $p_{*}(\pi_1(E))$ » est écrit au-dessus d'un mot biffé. Une
accolade ondulée dans la marge gauche court le long de l'énoncé.}

\textbf{Déf} : une arête de $E$ est un $1$-simplexe de $E$ muni d'un signe
$+$ ou $-$. Si $\sigma$ est un $1$-simplexe, on pose $\partial_1(+\sigma)
= \partial_1(\sigma)$, $\partial_0(+\sigma) = \partial_0(\sigma)$,
$\partial_1(-\sigma) = \partial_0(\sigma)$, $\partial_0(-\sigma) =
\partial_1(\sigma)$. Un chemin est une suite \add{d'arêtes} $\gamma =
(\gamma_1, \dots, \gamma_n)$ telle que $\partial_1(\gamma_{i+1}) =
\partial_0(\gamma_i)$. On pose $\partial_1(\gamma) = \partial_1(\gamma_1)$
et $\partial_0(\gamma) = \partial_0(\gamma_n)$. Un \struck{point}
\add{sommet} de base $x_0$ étant fixé dans $E$, un lacet est un chemin
$\gamma$ tel que $\partial_1(\gamma) = \partial_0(\gamma) = x_0$. 2 chemins
$\gamma$ et $\gamma'$ sont homotopes si leurs réalisations géométriques
sont homotopes modulo leurs extrémités. Cette relation d'équivalence est
compatible avec la loi de composition « mise bout à bout » des chemins, et
est engendrée par les suivantes :
\note{La page s'arrête sur les deux points ; la liste annoncée ne figure
pas sur la page suivante, qui s'ouvre sur un paragraphe biffé.}

\page{35}
\struck{Soit $E$ \ill{} un CW-complexe, $G$ un groupe discret opérant sur
$E$ par des transformations qui préservent la décomposition cellulaire,
$B = E/G$ \struck{l'ensemble} CW-complexe quotient de $E$}
\note{Trois lignes rayées de longs traits obliques.}

\textbf{Lemme 1} : $\forall x$ sommet de $E$, $\beta$ chemin de $B$ tel que
$\partial_1(\beta) = p(x)$, $\exists\, \gamma$ chemin de $E$ tel que
$\partial_1(\gamma) = x$ et $p(\gamma) = \beta$

\textbf{Dém} Il suffit de le voir si $\beta$ est une arête. Soit $\lambda$
une arête de $B$ telle que $p\lambda = \beta$. On a \struck{\ill{}}
$p(\partial_1 \lambda) = \partial_1(p\lambda) = \partial_1 \beta = p(x)$,
d'où $\exists\, g \in G$ \struck{\ill{}} $x = g.\partial_1 \lambda$, et
$\gamma = g\lambda$ répond à la question
\note{« $\lambda$ une arête de $B$ » : sic, pour une arête de $E$
au-dessus de $\beta$.}

\textbf{Lemme 2} : l'image de $p_{*} : \pi_1(E) \to \pi_1(B)$ est un
sous-groupe distingué de $\pi_1(B)$

\textbf{Dém} : Soit $\gamma$ un lacet de \struck{\ill{}} $E$, $\beta$ un
lacet de $B$, $\lambda$ un chemin de $E$ tel que $\partial_1(\lambda) =
x_0$ et $p_{*}(\lambda) = \beta$. On a $p(\partial_0 \lambda) = \partial_0
\beta = b_0 = p\, x_0$, d'où $\partial_0 \lambda = g x_0$, et le lacet
$(\lambda, g\gamma, -\lambda)$ se projette en $(\beta, p(\gamma), -\beta)$
ce qui démontre le lemme 2.

\textbf{Définition de $\chi : G \to \pi_1(B) / \operatorname{Im}
\pi_1(E)$ quand $E$ connexe.}

Pour tout $g \in G$, soit $\gamma$ un chemin de $E$ tel que
$\partial_1(\gamma) = x_0$, $\partial_0(\gamma) = g x_0$. \struck{\ill{}}
$p(\gamma)$ est un lacet de $B$, soit $\chi(g)$ sa classe dans $\pi_1(B) /
\operatorname{im} \pi_1(E)$. Elle ne dépend pas du chemin $\gamma$ choisi,
car, si $\gamma'$ est un autre chemin, $(\gamma, -\gamma')$ est un lacet de
$B$, et \struck{\ill{}} $p(\gamma)$ et $p(\gamma')$ ont même classe mod
l'image de $\pi_1(E)$.
\note{« $(\gamma, -\gamma')$ est un lacet de $B$ » : sic, pour un lacet de
$E$.}

\page{36}
Il résulte du lemme 1 que $\chi$ est surjective.

Montrons que $\chi$ est \struck{\ill{}} un homomorphisme : Soient $g_1, g_2
\in G$. Soit $\gamma_1$ un \struck{la} chemin tel que $\partial_1(\gamma_1)
= x_0$, $\partial_0(\gamma_1) = g_1 x_0$, idem \struck{\ill{}} $\gamma_2$
tel que $\partial_1(\gamma_2) = x_0$, $\partial_0(\gamma_2) = g_2 x_0$. On
a $\partial_1(g_1 \gamma_2) = g_1 x_0$ et $\partial_0(g_1 \gamma_2) = g_1
g_2 x_0$ : le chemin $\gamma = (\gamma_1, g_1 \gamma_2)$ vérifie
$\partial_1 \gamma = x_0$, $\partial_0 \gamma = g_1 g_2 x_0$. \struck{\ill{}
On a}
\[
p\gamma = (p\gamma_1, p g_1 \gamma_2) = (p\gamma_1, p\gamma_2) \struck{\ill{}}
\quad \text{d'où } \chi(g_1 g_2) = \chi(g_1)\chi(g_2)
\]
\note{« Soient $g_1, g_2 \in G$ » est ajouté en marge droite, au-dessus de
la ligne.}

Le noyau de $\chi$ contient évidemment le stabilisateur de $x_0$. Soit $x$
un sommet de $E$ et $s \in G$ tel que $sx = x$. Soit $\lambda$ un chemin de
$E$ tel que $\partial_1(\lambda) = x_0$, $\partial_0 \lambda = x$.
\struck{\ill{}} $\gamma = (\lambda, -s\lambda)$ est un chemin de $E$ tel
que $\partial_1(\gamma) = x_0$, $\partial_0 \gamma = s x_0$, d'où
$\chi(s) = \text{classe}\,(p(\lambda, -s\lambda)) = \text{classe}\,
(p(\lambda, -\lambda)) = $ élément neutre). Le noyau de $\chi$ contient donc
le sous-groupe $S$ \struck{\ill{}} engendré par les stabilisateurs des
\struck{\ill{}} sommets de $E$. Ce sous-groupe est distingué puisque
\struck{le} conjugué du stabilisateur de $x$ est le stabilisateur d'un
transformé de $x$.

On sait que $\operatorname{Ker} \chi \supset S$. Montrons qu'il ne le
contient pas strictement.

\textbf{Lemme 3} : $E/S$ est un revêtement \add{normal} de $E/G$, de groupe
$G/S$.

\page{38}
\note{La page 37 est blanche ; le texte se poursuit page 38.}
\struck{\ill{}}
\note{Une ligne gribouillée en tête de page.}

\textbf{Th 2} : Si $G$ est un groupe fini d'ordre $n$, l'image de $p_{*} :
H_{*}(E) \to H_{*}(B)$ contient tous les éléments \struck{\ill{}}
divisibles par $n$ dans $H_{*}(B)$.

\textbf{Dém} : Soit $\zeta$ \struck{\ill{}} $= \sum_{i \in I} a_i
\sigma_i$ un cycle de $B$, où $\sigma_i$ sont des simplexes. Pour tout
$\sigma_i$, soit $\tau_i$ simplexe de $E$ tel que $p\tau_i = \sigma_i$.
Alors
\[
\theta = \sum_{g \in G,\, i \in I} g\, a_i \tau_i
\]
est un cycle de $E$, et $p(\theta) = \struck{\ill{}}\ n\zeta \struck{\ill{}}$.
D'où le Th 2.

\textbf{Th 3} : Si \struck{\ill{}} $E$ est \struck{\ill{}} contractile,
$S =$ sous-groupe engendré par les stabilisateurs fini d'ordre $n$,
$\pi_i(B)$ est un groupe \struck{\ill{}} dont tous les éléments sont
d'ordre qui divise une puissance de $n$ pour $i > 1$
\note{« $S =$ sous-groupe engendré par les stabilisateurs » est écrit dans
l'interligne au-dessus de « fini d'ordre $n$ » ; « dont tous les éléments
sont d'ordre qui divise une puissance de $n$ pour $i > 1$ » descend dans la
marge gauche.}

\textbf{Dém} : \struck{\ill{}} D'après le Th.\ 1, $E/S$ est le revêtement
universel de $B$, d'où $\pi_i(B) = \pi_i(E/S)$. D'après le Th.\ 2,
\struck{\ill{}} les $H_i(E/S)$
\note{La démonstration s'arrête ici.}

\end{document}
